\section{From a question to a table}
\subsection{Define the object}

\begin{frame}{Before we begin}
  If the analysis needs one row for each country and year, what could produce two
  rows for the same country and year?

  \vspace{0.7em}
  \begin{block}{Write for two minutes}
    Give one possible cause and one check that could expose it.
  \end{block}

  \vfill
  We will return to these answers when we discuss joins.
\end{frame}

\begin{frame}{The goal for Lesson 2}
  \vfill
  \begin{center}
    {\LARGE\color{crimson}Turn one policy question into a checked analysis table.}

    \vspace{0.8em}
    The table should have a clear population, unit, key, period, variables, and
    record of the checks we performed.
  \end{center}
  \vfill
\end{frame}

\begin{frame}{What I want you to do by the end}
  \begin{enumerate}
    \item Translate a table contract into readable R.
    \item Use six data verbs for a stated reason.
    \item Predict whether rows, columns, values, or the unit should change.
    \item Test a country--year key before and after a join.
    \item Record how the table was built and checked.
    \item Ask an agent for one improvement that you can evaluate.
  \end{enumerate}
\end{frame}

\begin{frame}{Route for the lesson}
  \small
  \begin{tabularx}{\textwidth}{@{}>{\color{crimson}\bfseries}p{1.6cm}>{\bfseries}p{3.2cm}X@{}}
    \toprule
    Minutes & Work & Main question \\
    \midrule
    00--25 & Contract + inspect & What table do we need? \\
    25--60 & Select + filter & Which variables and observations remain? \\
    60--85 & Mutate + group & What value or unit did we create? \\
    85--103 & Join & Did any row multiply or fail to match? \\
    103--113 & Agent pass & Which bounded check can it improve? \\
    113--120 & Build record & Can another person inspect our table? \\
    \bottomrule
  \end{tabularx}
\end{frame}

\begin{frame}{The question comes before the verbs}
  Before we touch R, write:
  \begin{itemize}
    \item the population;
    \item the unit of observation;
    \item the key;
    \item the time period;
    \item the outcome;
    \item the comparison and grouping variables.
  \end{itemize}

  \vfill
  We call this the \alert{table contract}. It describes the object that the code
  should produce.
\end{frame}

\begin{frame}{The table contract is a map of the question}
  \begin{center}
  \begin{tikzpicture}[
    node distance=0.40cm,
    box/.style={draw=palegray,rounded corners=2pt,text width=2.30cm,
      minimum height=1.08cm,align=center,font=\small},
    arr/.style={-{Latex[length=1.8mm]},color=crimson,thick}
  ]
    \node[box] (question) {\textbf{Policy question}\\income and mortality};
    \node[box,right=of question] (contract) {\textbf{Table contract}\\population + unit + key};
    \node[box,right=of contract] (code) {\textbf{R pipeline}\\select + filter + mutate};
    \node[box,right=of code] (table) {\textbf{Checked table}\\one country--year row};
    \draw[arr] (question)--(contract);
    \draw[arr] (contract)--(code);
    \draw[arr] (code)--(table);
  \end{tikzpicture}
  \end{center}

  \vspace{0.7em}
  A map is selective. The contract records the choices needed to connect the
  source data to this question, and the build record shows whether the code
  followed that map.
\end{frame}

\begin{frame}{The health table contract}
  \footnotesize
  \renewcommand{\arraystretch}{1.16}
  \begin{tabularx}{\textwidth}{@{}>{\bfseries}p{2.4cm}p{5.2cm}X@{}}
    \toprule
    Field & Decision & Why it matters \\
    \midrule
    Population & Countries and economies with observed indicators & Defines who can appear. \\
    Unit & Country--year & Defines what one row means. \\
    Key & \texttt{iso3c + year} & Must identify every row once. \\
    Period & 2000--2022 & Defines the comparison window. \\
    Outcome & Under-five mortality & Names the quantity to describe. \\
    Comparison & GDP per capita, PPP & Names the relationship of interest. \\
    \bottomrule
  \end{tabularx}

  \vfill
  This is a descriptive exercise with World Development Indicators
  \citep{worldbankwdi}.
\end{frame}

\begin{frame}[fragile]{Read the frozen file}
\begin{lstlisting}[language=R]
library(tidyverse)

wdi <- read_csv(
  "data/derived/math-camp-wdi-2000-2022.csv",
  show_col_types = FALSE
)

glimpse(wdi)
\end{lstlisting}

  Keep the frozen file unchanged and build forward with code that can be run
  again \citep{wickham2023r4ds}.
\end{frame}

\begin{frame}{Look before you transform}
  Four small questions:
  \begin{enumerate}
    \item What does one row represent?
    \item Which columns identify that row?
    \item Are the identifiers unique?
    \item Which years and countries are present?
  \end{enumerate}

  \vfill
  \texttt{glimpse()}, \texttt{names()}, \texttt{count()}, and a short
  \texttt{summarise()} can answer these structural questions compactly.
\end{frame}

\begin{frame}[fragile]{Record the source structure}
\begin{lstlisting}[language=R]
wdi |>
  summarise(
    rows = n(),
    countries = n_distinct(iso3c),
    first_year = min(year),
    last_year = max(year),
    duplicate_keys = anyDuplicated(data.frame(iso3c, year))
  )
\end{lstlisting}

  A structure check tells us whether the source has the rows, variables, and
  keys that we expected.
\end{frame}
